% RAID 2026 Submission — Springer LNCS format
% 20 pages excluding references
\documentclass[runningheads]{llncs}

\usepackage{graphicx}
\usepackage{booktabs}
\usepackage{amsmath}
\usepackage{hyperref}
\usepackage{xcolor}
\usepackage{multirow}
\usepackage{algorithm}
\usepackage{algpseudocode}

% Anonymous submission
\renewcommand{\thefootnote}{}

\begin{document}

\title{Evolved Feature Selection in Weightless Neural Networks\\for Network Intrusion Detection}

% Anonymous for review
\author{Anonymous}
\institute{Anonymous Institution}

\maketitle

% ============================================================================
\begin{abstract}
Weightless Neural Networks (WNNs) based on Random Access Memory (RAM) neurons
offer constant-time inference with sub-kilobyte model sizes, making them
attractive for edge-deployed intrusion detection systems. However, their
classification performance depends critically on \emph{connectivity}---which
input bits each neuron observes---a problem traditionally solved by random
initialization. We present an evolutionary optimization framework that
discovers optimal neuron connectivity through genetic algorithm and tabu
search, effectively performing automated feature selection intrinsic to the
architecture. We evaluate on three standard IDS benchmarks (UNSW-NB15,
CICIDS2017, CIC-IoT-2023) using both temporal and random splits, with a
rigorous eight-mode threshold calibration protocol that exposes the sensitivity
of binary classifiers to decision boundary selection. Across 100+ independent
runs per configuration, our approach achieves [TODO: headline numbers]
competitive with Random Forest and gradient-boosted models at 50,000$\times$
smaller model size ($<$1\,KB vs.\ $\sim$50\,MB), with inference requiring only
a single memory lookup per neuron.

\keywords{Intrusion Detection \and Weightless Neural Networks \and
Evolutionary Optimization \and Feature Selection \and FPGA}
\end{abstract}

% ============================================================================
\section{Introduction}
\label{sec:intro}

% Hook: IDS at edge needs speed + small models
Network intrusion detection systems (NIDS) face a fundamental tension between
detection accuracy and deployment constraints. Deep learning models achieve
state-of-the-art accuracy but require megabytes of parameters and GPU
inference, making them impractical for line-rate packet classification on edge
devices. Conversely, lightweight rule-based systems operate at wire speed but
lack the adaptability to detect novel attack patterns.

% Gap: WNNs exist for IDS but limited evaluation
Weightless Neural Networks (WNNs)~\cite{aleksander1984,ludermir1999} offer a
compelling middle ground: RAM-based neurons perform classification via direct
memory lookup in $O(1)$ time, with model sizes measured in bytes rather than
megabytes. Recent work~\cite{susskind2023fwiw} demonstrated WNN-based IDS
achieving 98\% accuracy on UNSW-NB15, but only under random train/test splits
that inflate performance through data leakage. The critical question of
connectivity optimization---which input bits each neuron should observe---was
left to random initialization.

% Our contribution
We address both limitations. First, we introduce an evolutionary optimization
framework (genetic algorithm + tabu search) that discovers optimal neuron
connectivity, effectively performing automated feature selection intrinsic to
the WNN architecture. Second, we evaluate rigorously on three datasets using
temporal splits that simulate realistic deployment, with a comprehensive
threshold calibration study that quantifies the gap between train-time and
deployment-time performance.

% Contributions list
Our contributions are:
\begin{enumerate}
\item \textbf{Evolved connectivity as feature selection.} GA/TS optimization
  of neuron connections provides the largest single improvement in
  classification performance, discovering which input features matter
  without external feature selection.
\item \textbf{Sub-kilobyte models with nanosecond inference potential.}
  Our evolved WNN achieves competitive F1 with Random Forest and XGBoost
  at 50,000$\times$ smaller model size ($<$1\,KB vs.\ $\sim$50\,MB),
  with $O(1)$ lookup-based inference suitable for FPGA deployment.
\item \textbf{Eight-mode threshold calibration study.} We identify threshold
  selection as a critical but under-reported factor in IDS evaluation,
  showing [TODO: X]\% F1 gap between oracle and train-calibrated thresholds.
\item \textbf{Three-dataset evaluation.} Results on UNSW-NB15, CICIDS2017,
  and CIC-IoT-2023 demonstrate architecture generalizability across
  different network environments and attack types.
\end{enumerate}

% ============================================================================
\section{Background}
\label{sec:background}

\subsection{RAM Neurons and Weightless Neural Networks}

% RAM neuron basics
Unlike conventional neural networks that learn continuous weights through
gradient descent, Weightless Neural Networks (WNNs)~\cite{aleksander1984}
use Random Access Memory (RAM) neurons that perform classification via
direct memory lookup. A RAM neuron consists of $2^n$ memory cells, where
$n$ is the number of observed input bits---the neuron's \emph{address width}.
Given a binary input vector of length $d$, the neuron selects $n$ bits
through a fixed \emph{connectivity map} $C = \{c_1, \ldots, c_n\}$ where
$c_i \in \{0, \ldots, d{-}1\}$, concatenates them to form an $n$-bit
address, and returns the stored value at that address. Training writes
target labels to addressed cells; inference reads them in $O(1)$ time
with no multiply-accumulate operations.

% Generalization through partial connectivity
The fundamental mechanism enabling WNN generalization is
\emph{partial connectivity}. A fully-connected neuron ($n = d$) memorizes
every training example without generalizing, since each input maps to a
unique address. With partial connectivity ($n \ll d$), many distinct inputs
share the same address---those that agree on the $n$ selected bit
positions---causing the neuron to learn a \emph{feature pattern} over its
observed bits. An ensemble of $N$ neurons with diverse connectivity maps
provides $N$ complementary views of the input, and their aggregated scores
produce a classification decision.

% QUAD_WEIGHTED memory
\textbf{Quad-weighted memory.}
Traditional RAM neurons store binary values (TRUE/FALSE). We extend this
with \emph{quad-weighted} (QW) memory using four states that encode
graduated confidence: FALSE ($0.0$), WEAK\_FALSE ($0.25$), WEAK\_TRUE
($0.75$), and TRUE ($1.0$). Each cell is stored in 2 bits. When a cell
is written with a conflicting label, it transitions to the corresponding
weak state rather than being overwritten, preserving partial information
from both classes. The neuron's output score---the mean of its $N$
neurons' cell values---thus provides a smoother probability estimate than
binary voting.

\subsection{Thermometer Encoding}

WNNs require binary inputs, but network flow features (packet counts,
byte rates, inter-arrival times) are continuous. We convert each feature
to a binary vector using \emph{thermometer encoding}~\cite{susskind2022bthowen}.
Given a feature value $x$ and $k$ thresholds $\theta_1 < \cdots < \theta_k$,
the encoding produces $k$ bits: bit $i$ is 1 if $x \geq \theta_i$, else 0.
This preserves ordinal relationships---similar values differ by few bits,
enabling RAM neurons to generalize across nearby feature values.

We compute thresholds using the \emph{distributive} (quantile-based) method,
which places thresholds at equally-spaced quantiles of the training
distribution. This is robust to skewed distributions common in network
traffic (e.g., packet sizes follow a heavy-tailed distribution). With $k$
bits per feature and a feature set $F$, the total input width is
$d = k \times |F|$ bits.

\subsection{Evaluation Datasets}

We evaluate on three standard IDS benchmarks, each with distinct
characteristics:

\textbf{UNSW-NB15}~\cite{moustafa2015unswnb15} contains network flow records
from the Australian Centre for Cyber Security, with 44 features and 10 attack
categories. The standard temporal split (175K train / 82K test) separates
training and testing periods, preventing data leakage. We also evaluate on a
deduplicated random split (1.4M / 158K) for comparison with prior
work~\cite{susskind2023fwiw}.

\textbf{CICIDS2017}~\cite{sharafaldin2018cicids} comprises five days of
network traffic captured at the Canadian Institute for Cybersecurity. We
define a temporal split using Monday--Thursday for training (2.1M flows) and
Friday for testing (703K flows), requiring the model to generalize to
previously unseen attack types (DDoS, Botnet, PortScan). This split has
not been previously used in the literature---most work uses random splits
that achieve $>$99\% accuracy.

\textbf{CIC-IoT-2023}~\cite{neto2023ciciot} captures IoT traffic from 105
heterogeneous devices with 33 attack types across 7 classes. With 46.7M
records heavily skewed toward attacks (97.6\%), we subsample to 1.3M records
for a balanced benchmark. Only a random split is feasible (no temporal
ordering in the source data).

% ============================================================================
\section{Related Work}
\label{sec:related}

\subsection{WNN-Based Intrusion Detection}

The WiSARD~\cite{ludermir1999} architecture pioneered RAM-based pattern
recognition, and several works have applied WNNs to network security.
Santiago et al.~\cite{santiago2020} demonstrated WNN-based anomaly detection
but evaluated only on random splits with random connectivity.

The most directly related work is FWIW~\cite{susskind2023fwiw}, which
achieves 98\% accuracy on UNSW-NB15 with a 272-byte model deployed on FPGA.
FWIW uses Bloom filter-based neurons with H3 hashing and backpropagation
training. However, three limitations remain: (1)~evaluation uses only random
splits, inflating accuracy through train/test data leakage; (2)~neuron
connectivity is randomly initialized, leaving significant accuracy on the
table; and (3)~only a single threshold mode is reported, obscuring the
sensitivity of binary classification to decision boundary selection.

BTHOWeN~\cite{susskind2022bthowen} introduced thermometer encoding with
learned thresholds for WNN classification, achieving strong results on
tabular benchmarks. We adopt their encoding scheme but extend it with
distributive (quantile-based) thresholds and evolutionary connectivity
optimization.

\subsection{Machine Learning for IDS}

Traditional ML approaches---Random Forest, XGBoost, SVM---achieve
85--90\% F1 on UNSW-NB15's temporal split~\cite{zoghi2024} but require
models of 10--50\,MB. Deep learning approaches (BiLSTM~\cite{mdpi2024enhanced},
CNN-LSTM) push accuracy higher at the cost of GPU inference requirements.

A critical methodological issue pervades the IDS literature: most published
results use random train/test splits, reporting 97--99\% accuracy that does
not reflect deployment performance. Zoghi and Serpen~\cite{zoghi2024} showed
that the same models achieving 99\% on random splits drop to 87\% on the
standard temporal split, a finding consistent with our evaluation. For
CICIDS2017, the situation is similar---near-perfect random-split results
contrast with the absence of temporal evaluation in the literature. Known
labeling errors~\cite{engelen2021cicids} further complicate comparisons.

\subsection{FPGA-Based Network Security}

FPGA acceleration for network security spans from packet-level inspection
(Pigasus~\cite{zhao2020pigasus}) to ML inference
(LogicNets~\cite{umuroglu2020logicnets}, PolyLUT~\cite{andronic2023polylut}).
WNNs are particularly well-suited for FPGA deployment because RAM neurons
map directly to FPGA lookup tables (LUTs) and block RAM (BRAM), requiring
no multiply-accumulate units. FWIW~\cite{susskind2023fwiw} demonstrated
this with sub-microsecond inference on a Virtex UltraScale+. Our work
extends this direction with evolved connectivity that can further reduce
resource utilization by discovering sparser, more efficient neuron
configurations.

% ============================================================================
\section{Architecture}
\label{sec:architecture}

\subsection{Single-Cluster Binary Discriminator}

% Architecture overview
Our IDS architecture uses a \emph{single-cluster discriminator}: $N$ RAM
neurons with $b$ address bits each, all observing the same thermometer-encoded
input. Each neuron independently classifies the input as normal (0) or attack
(1) by looking up its addressed memory cell. The cluster's output is the mean
of neuron scores, thresholded to produce a binary decision.

% Why single-cluster works
This deliberately simple architecture makes the connectivity optimization
problem tractable: the GA/TS only needs to optimize $N \times b$ connection
indices rather than per-class architectures.

\subsection{Phased Evolutionary Optimization}

% Multi-phase pipeline
We optimize the discriminator through a phased pipeline:

\begin{enumerate}
\item \textbf{Grid search.} Exhaustive evaluation of $(N, b)$ configurations,
  ranking by a weighted harmonic mean of CE, F1, and FPR.
\item \textbf{GA neurons.} Genetic algorithm optimizing neuron count per
  cluster while preserving learned connections.
\item \textbf{GA bits.} Optimizing address width (bits per neuron).
\item \textbf{GA connections.} The critical phase---optimizing \emph{which}
  input bits each neuron observes. This is where the architecture learns
  its feature selection.
\end{enumerate}

% Fitness function
The fitness function uses weighted harmonic mean of ranks across four metrics:
\begin{equation}
\text{fitness} = \frac{w_{CE} + w_{F1} + w_{FPR}}{%
  \frac{w_{CE}}{r_{CE}} + \frac{w_{F1}}{r_{F1}} + \frac{w_{FPR}}{r_{FPR}}}
\end{equation}
where $r_m$ is the genome's rank for metric $m$ within the population.
Default weights: $w_{CE}=0.4$, $w_{F1}=0.3$, $w_{FPR}=0.3$.

\subsection{Rust+Metal GPU Acceleration}

% System contribution
[TODO: Brief description of the Rust accelerator with Metal GPU support.
Key point: makes population-based optimization practical, enabling 100+
independent runs for statistical analysis.]

% ============================================================================
\section{Evaluation Methodology}
\label{sec:methodology}

\subsection{Temporal vs.\ Random Splits}

% Why this matters
[TODO: Explain temporal vs random splits, data leakage in random splits,
why temporal is the honest evaluation. UNSW-NB15: 175K/82K temporal.
CICIDS2017: Mon-Thu/Friday temporal. CIC-IoT-2023: random only.]

\subsection{Eight-Mode Threshold Calibration Protocol}

% The threshold problem
Binary IDS classifiers produce a continuous score that must be thresholded.
The choice of threshold dramatically affects F1 and FPR, yet most papers
report only a single threshold mode (typically train-calibrated or fixed 0.5).

We evaluate each genome under eight threshold modes:
\begin{enumerate}
\item \textbf{Train-cal}: Threshold optimized on training scores
\item \textbf{Test-cal}: Threshold swept on held-out validation data
\item \textbf{Fixed 0.5}: No calibration (distribution-agnostic baseline)
\item \textbf{Platt scaling}: Sigmoid calibration on held-out scores
\item \textbf{Beta calibration}: Logit-based calibration
\item \textbf{Empirical}: Histogram-based threshold
\item \textbf{Empirical cumulative}: CDF-based threshold
\item \textbf{Oracle (val-cal)}: Threshold swept on test data (upper bound)
\end{enumerate}

\subsection{Statistical Protocol}

% 34+ runs per config
Each configuration is run [TODO: 34--100] times with different random seeds.
We report mean $\pm$ standard deviation and 95\% confidence intervals.
All results are from \texttt{validation\_summaries} (per-flow honest
evaluation), not cherry-picked leaderboard entries.

% ============================================================================
\section{Experimental Results}
\label{sec:results}

\subsection{UNSW-NB15 Results}

% Main table: 34 flows, all threshold modes
[TODO: Table from ids\_results.md — top20-CE4F3R3 (34 flows)]

\subsection{CICIDS2017 Results}

% Temporal split results
[TODO: Results from 34 CICIDS2017 temporal flows once completed]

\subsection{CIC-IoT-2023 Results}

% Random split results (stretch goal)
[TODO: Results if completed before deadline]

\subsection{Phase Progression Analysis}

% Which phases help/hurt
[TODO: Table showing grid_search → ga_neurons → ga_bits → ga_connections
progression for each architecture]

\subsection{Comparison with Baselines}

% The money table
[TODO: Comparison table — RF, XGBoost, BiLSTM, FWIW vs Our WNN.
Separate temporal and random split results.]

\subsection{Model Size Analysis}

% Sub-1KB claim
[TODO: Compute actual model sizes. N neurons × b bits × 2 bits/cell ÷ 8 = bytes.
Compare with RF (~50MB), DNN (~5MB), FWIW (272B).]

% ============================================================================
\section{Discussion}
\label{sec:discussion}

\subsection{Connectivity as Feature Selection}

[TODO: The GA connection optimization phase provides the largest accuracy
improvement. Analysis of which features the evolved neurons select vs.
RF importance ranking.]

\subsection{Threshold Sensitivity}

[TODO: The ~5% F1 gap between oracle and train-cal represents the
calibration challenge. Platt/beta calibration partially closes this gap.
Implications for deployment.]

\subsection{FPGA Deployment Potential}

% Not implemented, but analyzed
[TODO: BRAM requirements for our genomes. A neuron with $b=32$ address
bits needs $2^{32} \times 2$ bits = 1\,GB --- too large for direct mapping.
But with sparse memory (most cells empty), only non-empty entries need
storage. Typical occupancy analysis. Comparison with FWIW FPGA and
LogicNets.]

\subsection{Limitations}

[TODO: Binary classification only (multi-class as future work).
Threshold sensitivity requires calibration data.
Temporal split not available for CIC-IoT-2023.]

% ============================================================================
\section{Related Work}
\label{sec:related}

\subsection{WNN for Intrusion Detection}

[TODO: FWIW~\cite{susskind2023fwiw}, Santiago et al., BTHOWeN~\cite{susskind2022bthowen}]

\subsection{ML-Based IDS}

[TODO: Zoghi 2024, Kasongo 2020, survey of RF/XGBoost/DL approaches on
UNSW-NB15 and CICIDS2017]

\subsection{FPGA-Based IDS}

[TODO: LogicNets, PolyLUT, Pigasus. Hardware acceleration landscape.]

% ============================================================================
\section{Conclusion}
\label{sec:conclusion}

[TODO: Summary of contributions. Key result highlight. Future work:
multi-class hierarchical classification, online learning for threshold
adaptation, cross-dataset transfer, FPGA implementation.]

% ============================================================================
\bibliographystyle{splncs04}
\bibliography{references}

\end{document}
